%% ============================================================
%% 片段：s06-local   对应原文 PDF 第 130--145 页
%% 本文件由 gw/convert.py 从 arXiv 源码（hep-th/0612073v2）自动生成：
%%   * 公式/编号/引用/图 已经锁定为原文内容，请勿改动任何数学内容与编号；
%%   * 只需把英文正文译为中文，并在每节开头加“本节重点”、关键处加译注。
%% ============================================================


\section{局部模型与弦论实现}

%% ==== 原文 §6：PDF 第 130 页（印刷第 129 页）====
\label{sec:6}

\begin{keytakeaways}
  \begin{itemize}[leftmargin=1.4em,itemsep=2pt,topsep=2pt]
    \item \textbf{物理图像}：面算符可用“局部模型”实现，即四维 ${\mathcal N}=4$ 超杨–米尔斯理论与支撑在二维子流形上的超对称西格玛模型的耦合；后者的靶空间为 $T^*(G/\mathbb{L})$，其超 Kähler 度量来自余伴随轨道。局部模型只适用于 $\MH$ 的局部奇点附近，且不保持 $S$-对偶。
    \item \textbf{数学工具}：$G_\C=SL(2,\C)$ 时用线性西格玛模型：$\mathbb{H}^2$ 的 $U(1)$ 超 Kähler 商给出 $A_1$ 奇点的超 Kähler 消解，Fayet–Iliopoulos 参数对应 $(\alpha,\beta,\gamma)$，$\theta$ 角对应 $\eta$；$O(2)$ 的非连通部分给出外尔群作用。
    \item \textbf{关键结论}：局部模型的单值群是轨形 $\tilde\C^*/\Z_2$ 的基本群；保持离散电通量 ${\bf e}_0$ 的子群恰为 $SL(2,\C)$ 的仿射辫群，这与第~\ref{sec:4.5}~节的结论一致。
    \item \textbf{瞬子图像}：局部模型在 $(\alpha,\beta,\gamma,\eta)=0$ 处的奇点来自小作用量瞬子级数在 $q=1$ 的极点；在四维中可用消没定理证明，只有 $\alpha/\gamma=(t-t^{-1})/2$、$\beta=0$ 时才存在这类瞬子。
    \item \textbf{弦论实现}：相交的 $D3$/$D3'$ 膜给出 $\mathbb{H}^N/\neg/\neg/U(1)$ 型线性西格玛模型，但破坏 $S$-对偶；改用 AdS/CFT（$AdS_5\times S^5$ 中的 $AdS_3\times S^1$）或 M 理论（$AdS_7\times S^4$ 中的 $M2$-膜）即可保持 $S$-对偶并获得面算符。
  \end{itemize}
\end{keytakeaways}


\subsection{概述}
%% ==== 原文 §6.1：PDF 第 130 页（印刷第 129 页）====
\label{sec:6.1}

本文的最后一节讨论若干前文主题的另类进路。我们先作一概览。

与其只靠假定场具有熟悉的奇点
% 式 (6.1)
\setcounter{equation}{0}
\begin{equation}\label{eq:6.1}
  \begin{aligned}A & =  \alpha \,d\theta \\
\phi & = \beta\,{dr\over r}-\gamma\,d\theta\end{aligned}
\end{equation}
（沿余维二的曲面 $D$）来描述面算符，人们自然会想用另一方式描述它：引入支撑在 $D$ 上的新自由度，并把它们与体（bulk）中 ${\mathcal N}=4$ 超杨–米尔斯理论的规范场耦合起来。例如，可把二维中具有 $(4,4)$ 超对称的超多重态（hypermultiplet）或向量多重态（vector multiplet）耦合到体中的 ${\mathcal N}=4$ 超杨–米尔斯理论。这样给出的理论，在某种程度上类似我们在本文中“手工”描述过的理论。

物理文献中当然已有许多关于共形场论与各种缺陷（defect）耦合的分析；所谓缺陷，即支撑在正余维子流形上的场。凝聚态物理中的 Kondo 模型（Kondo model）是一个经典例子。与当前关切更接近的是：定域在余维一的超多重态已通过膜构造（brane construction）与 ${\mathcal N}=4$ 超杨–米尔斯理论耦合，例如见 \cite{ref78,ref79}；而与我们最相关的余维二超多重态则已在 \cite{ref18} 中讨论过。

{}从第~\ref{sec:3.8}~节，我们其实可以看出什么样的杂质模型（impurity model）能重现我们所考虑的面算符的某些性质。本文的重要主题之一是希格斯丛模空间 $\MH$ 在三元组 $(\alpha,\beta,\gamma)$ 变为非正则时的“局部奇点”（local singularity）。这种局部奇点的超 Kähler 消解（hyper-Kähler resolution），可方便地用某些复余伴随轨道（coadjoint orbit）上的超 Kähler 度量~\cite{ref34} 来描述；等价地（在另一个复结构下），可用各列维子群 $\mathbb{L}\subset G$ 对应的 $X=T^*(G/\mathbb{L})$ 上的超 Kähler 度量来描述。因此，本文探讨的某些主题实际上可以用一个复合系统来表达：四维时空上的 ${\mathcal N}=4$ 超杨–米尔斯理论，加上支撑在时空的某个二维子流形上的超对称西格玛模型。我们称这个西格玛模型为局部模型（local model）。

尽管这一进路很有用，但据我们所知它有一个根本性的局限。本文发展的理论依赖于参数\footnote{为简洁起见，这里描述的是列维子群为 $\mathbb{L}=\TT$ 时的参数。一般地，所有参数都要求是 $\mathbb{L}$-不变的。} $(\alpha,\beta,\gamma,\eta)
\in \TT\times \mathfrak t\times\mathfrak t\times {}^L\neg\TT$。西格玛模型依赖于大体类似的参数 $(\alpha,\beta,\gamma,\eta)$：其中 $(\alpha,\beta,\gamma)$ 是进入 $X$ 上超 Kähler 度量的几何参数，$\eta$ 则包含西格玛模型的 $\theta$ 角。参数 $\eta$ 取值于 $^L\neg\TT$，与规范场论中一样（见式~\eqref{eq:3.44}）；同样，两种情形下 $\beta$ 与 $\gamma$ 都取 $\mathfrak t$ 值。但 $\alpha$ 在 $\MH$ 的完整几何中取 $\TT$ 值，而在 $T^*(G/\mathbb{L})$ 的超 Kähler 几何中取 $\mathfrak t$ 值（并通过 $SO(3)$ 对称性与 $\beta,\gamma$ 相联系），如我们在第~\ref{sec:3.8}~节回顾过的。因此，完整模型与局部模型的参数并不完全一致。

这与“局部模型刻画局部奇点附近行为”的说法并不矛盾，因为对这一目的而言，$\alpha$ 变大时的行为并不重要。然而，这显然意味着局部模型不可能具有 $S$-对偶。既然 $S$-对偶交换 $\alpha$ 与 $\eta$，它在靶为 $T^*(G/\mathbb{L})$ 的西格玛模型中就很难成立：在该模型中 $\eta$ 是角变量，而 $\alpha$ 取值于线性空间。因此，若把 ${\mathcal N}=4$ 理论耦合到这个局部模型，就不能指望保持 $S$-对偶。

这或许显得奇怪，但换一个观点也许会使该结论显得更自然。空间 $T^*(G/\mathbb{L})$ 不是环面纤维化，也不具有镜像对称（mirror symmetry），至少不具有 $G$-不变的镜像对称。（若西格玛模型的镜像对称不是 $G$-不变的，那么当它耦合到规范群为 $G$ 的四维规范场论时也帮不上多少忙。）因此，指望在把 ${\mathcal N}=4$ 超杨–米尔斯理论耦合到局部模型时仍保持 $S$-对偶，也许是不现实的。

尽管如此，正如我们强调过的，我们确实相信局部模型有助于研究 $\MH$ 的局部奇点附近（即非正则三元组 $(\alpha,\beta,\gamma)$ 附近）的行为。例如，考虑仿射辫群（affine braid group）$\BAff$ 在 $\Psi=\infty$ 处 $B$-模型的膜上的作用。在此作用中，$\alpha$ 与 $\gamma$ 保持固定（最有趣的情形是它们接近零），而研究的是关于 $\beta$ 与 $\eta$ 的单值性。$\alpha$ 取 $\mathfrak t$ 值还是 $\TT$ 值并不重要；$\BAff$ 作用于局部模型的 $B$-膜，其理由与它作用于 $\MH$ 上的膜相同。事实上，\cite{ref29} 中对 $\BAff$ 作用的分析，从物理观点看等价于研究局部模型 $B$-膜上的单值性作用。对于 $\MH$ 的 $B$-膜上的 $\BAff$ 作用，这似乎是一个好的局部模型。

相反，局部模型不是研究 $\Psi=0$ 处 $A$-膜上 $\BAff$ 作用的好模型。这涉及对 $(\alpha,\beta)$ 这对参数的单值性作用，而一旦 $\alpha$ 取 $\mathfrak t$ 值而非 $\TT$ 值，这一单值性作用就丢失了。据我们所知，要研究 $\MH$ 的 $A$-膜上的单值性作用，必须使用靶为 $\MH$ 的完整西格玛模型（或完整的四维规范场论）。这里有一个略带讽刺的角色反转：$A$-模型的单值性作用涉及经典辛同胚，而 $B$-模型的作用则高度非经典；但 $B$-模型的单值性作用可以在局部模型中研究，$A$-模型的却不能。

\bigskip\noindent{\it 本节内容}

现在我们可以简要概述本节内容。

正如已讨论的，尽管有局限，局部模型确有其用处。随之而来的问题是：是否存在关于局部模型的某种物理进路，使我们能比已获得的了解更多？对一般规范群我们尚无好的答案，但对 $G_\C=SL(2,\C)$ 的情形，局部模型有一个有用的描述。这将是第~\ref{sec:6.2}~节的主题。

第~\ref{sec:6.2}~节的分析将揭示一个事实的重要性：当三元组 $(\alpha,\beta,\gamma)$ 变为非正则时，（局部模型的某些超对称瞬子的）瞬子作用量消失。这一事实其实继承自四维规范场论中的类似事实，我们将在第~\ref{sec:6.3}~节说明。

最后，仅仅说局部模型缺少对偶对称性并不令人满意。我们能改用什么类似、但确实保持对偶对称性的做法？事实上，弦理论~\cite{ref18} 已在一个有趣的特别情形（$G_\C=SL(N,\C)$ 配以极小余伴随轨道）给出了这一问题的答案。我们在第~\ref{sec:6.4}~节说明。

\subsection{\texorpdfstring{$G_\C=SL(2,\C)$ 的线性西格玛模型}{G_C = SL(2,C) 的线性西格玛模型}}
 %% ==== 原文 §6.2：PDF 第 132 页（印刷第 131 页）====
\label{sec:6.2}

要利用局部模型——即靶为 $T^*(G/\mathbb{L})$ 的超对称西格玛模型——来学到四维规范场论中不明显的东西，我们需要一种不同的研究方式。我们没有对所有群一致适用的普遍方法，但将解释 $G=SU(2)$ 的一种有趣进路。这里我们依赖的事实是：唯一相关的 $\mathbb{L}$ 选取是 $\mathbb{L}=U(1)$，且 $T^*(SU(2)/U(1))=T^*\mathbb{CP}^1$ 是 $A_1$ 奇点的消解。因此，它可以构造为有限维线性空间的超 Kähler 商（hyper-Kähler quotient）~\cite{ref62}。这在膜物理中有许多应用~\cite{ref80}，且与我们的目的相关，因为可用它构造线性西格玛模型（linear sigma model）~\cite{ref81}，从而简单地理解靶为 $T^*\mathbb{CP}^1$ 的超对称西格玛模型的真空结构。

所有典型李群的复余伴随轨道同样可以构造为有限维线性空间的超 Kähler 商~\cite{ref82}。这些更一般的构造可以作为构造适用于任意典型群、与面算符相关的线性西格玛模型的起点。然而，$SU(2)$ 情形有一个重要优势：外尔群的作用在线性西格玛模型中一目了然。对更高秩的 $G$ 则并非如此，这就是我们在本节分析中只考虑 $SU(2)$ 的原因。

\begin{yizhu}[物理图像补充]这里的构造就是 Kronheimer 的 ALE 空间构造~\cite{ref62}：在 $\mathbb{H}^2\cong\R^8$ 上按三个矩映射分量 $(\alpha,\beta,\gamma)$ 取 $U(1)$ 超 Kähler 商，参数全为零时得到 $A_1$ 奇点 $\R^4/\Z_2$，否则得到它的超 Kähler 消解。一般典型群的复余伴随轨道同样可实现为有限维线性空间的超 Kähler 商~\cite{ref82}，其超 Kähler 度量由 Kronheimer 给出~\cite{ref34}。这正是本文局部模型取 $X=T^*(G/\mathbb{L})$ 的依据：在某个复结构下 $X$ 是消解，在另一个复结构下同构于相应的复余伴随轨道；$SU(2)$ 情形即 $T^*\mathbb{CP}^1$ 与 $A_1$ 消解。由于只有 $SU(2)$ 的外尔群作用能在线性模型中清楚看到（见下文 $O(2)$ 的非连通部分），本节才限定 $SU(2)$。\end{yizhu}

为把 $A_1$ 奇点实现为超 Kähler 商，我们从 $\mathbb{H}^2\cong \R^8$ 出发，视其为平坦的超 Kähler 流形。它承载 $Sp(2)$ 的作用；在我们的记号中，$Sp(2)$ 是 $2\times 2$ 四元数酉矩阵构成的群。两类四元数矩阵将起重要作用。$Sp(2)$ 包含由单位阵的四元数倍数构成的群 $SU(2)$
% 式 (6.2)
\setcounter{equation}{1}
\begin{equation}\label{eq:6.2}
  \left(\begin{matrix}v&0\\ 0&v\end{matrix}\right),
\end{equation}
其中 $v$ 是模为 $1$ 的四元数（quaternion）；它还包含由元素全为实数的酉四元数矩阵构成的群 $O(2)$：
% 式 (6.3)
\setcounter{equation}{2}
\begin{equation}\label{eq:6.3}
  \left(\begin{matrix}a&b\\ c&d\end{matrix}\right).
\end{equation}

这两个群彼此交换，但整体结构并非直积，因为 $SU(2)$ 与 $O(2)$ 含有共同的中心元素 $-1$。因此整体结构是
% 式 (6.4)
\setcounter{equation}{3}
\begin{equation}\label{eq:6.4}
  \left(SU(2)\times O(2)\right)/\Z_2\subset Sp(2).
\end{equation}

\begin{yizhu}[易错点]两个因子群并非直积：$SU(2)$ 与 $O(2)$ 的交恰为共同中心 $\Z_2=\{\pm 1\}$，所以式 \eqref{eq:6.4} 是商 $(SU(2)\times O(2))/\Z_2$。这里的 $Sp(2)$ 指 $2\times 2$ 四元数酉矩阵群，实维数 $10$，同构于 $Spin(5)$；而 $SU(2)\times O(2)$ 的连通部分只有 $3+1=4$ 维。读式 \eqref{eq:6.4} 时不要误认成直积：正是这个公共 $\Z_2$ 使 $O(2)$ 的非连通元素 $w$ 能同时充当 $U(1)$ 的外自同构与 $SU(2)$ 外尔群的非平凡元。\end{yizhu}


$O(2)$ 的连通分支是 $U(1)$。令 $\vec \mu=(\mu_1,\mu_2,\mu_3)$ 为这个 $U(1)$ 作用的超 Kähler 矩映射。引入实参数 $(\alpha,\beta,\gamma)$，并把 $X(\alpha,\beta,\gamma)$ 定义为在 $\vec\mu=(\alpha,\beta,\gamma)$ 处取值的超 Kähler 商 $\mathbb{H}^2/\neg/\neg/U(1))$。换言之，$X(\alpha,\beta,\gamma)$ 由令 $\vec\mu=(\alpha,\beta,\gamma)$ 并按 $U(1)$ 取商来定义。当 $\alpha=\beta=\gamma=0$ 时，这给出 $A_1$ 奇点 $\R^4/\Z_2$；一般地，它给出 $A_1$ 奇点的超 Kähler 消解。

\begin{yizhu}[易错点]三斜杠商 $\mathbb{H}^2/\neg/\neg/U(1)$ 表示超 Kähler 商：固定超 Kähler 矩映射 $\vec\mu$ 的三个分量后按 $U(1)$ 取商，而非通常的复商或 GIT 商；它在稳定点处与辛商一致，但保留超 Kähler 结构。三个 Fayet–Iliopoulos 参数就是矩映射的三个分量，$SO(3)$ 对称性旋转它们。注意原文 PDF 第 133 页（印刷第 132 页）在此处多排了一个右括号（$\mathbb{H}^2/\neg/\neg/U(1))$），按定义应读作 $\mathbb{H}^2/\neg/\neg/U(1)$；译文照录原文，未改公式。\end{yizhu}

{}从物理观点看，我们在二维流形 $D$ 上构造一个超对称线性西格玛模型，其场为一个 $U(1)$ 向量多重态和一对超多重态 $H_1,\, H_2$（其玻色分量参数化 $\mathbb{H}^2$，并按上述方式受 $U(1)$ 作用）。参数 $(\alpha,\beta,\gamma)$ 是向量多重态的 $D$-项，即 Fayet–Iliopoulos 参数。经典真空的模空间恰是上一段描述的超 Kähler 商 $X(\alpha,\beta,\gamma)$，但有一点需注意：当 $\alpha=\beta=\gamma=0$ 时，该模空间还有第二个分支（second branch），如 \cite{ref81} 所述。

线性西格玛模型具有旋转参数 $(\alpha,\beta,\gamma)$ 的 $SO(3)$ 对称性。其来源如下。若把 $\mathbb{H}^2$ 的元素看作四元数列 $\left(\begin{matrix}u_1\\ u_2\end{matrix}\right)$，则 $Sp(2)$ 的左作用与 $Sp(1)$ 的右作用交换。这个 $Sp(1)$ 的作用是旋转超 Kähler 流形 $\mathbb{H}^2$ 的复结构，并同样旋转参数 $(\alpha,\beta,\gamma)$。

除 $(\alpha,\beta,\gamma)$ 外，量子理论还依赖于 $U(1)$ 向量多重态的 $\theta$ 角，照例我们将其记作 $\eta$。显然 $\eta$ 是角变量。它以如下形式的因子出现在路径积分中
% 式 (6.5)
\setcounter{equation}{4}
\begin{equation}\label{eq:6.5}
  \exp\left(i\eta\int_DF\right).
\end{equation}
这里 $F$ 是 $U(1)$ 曲率，我们用一个实二维形式表示它。这意味着，在模型的通常形式中，存在对称性
% 式 (6.6)
\setcounter{equation}{5}
\begin{equation}\label{eq:6.6}
  \eta\to \eta+1.
\end{equation}
正如我们将在后文解释的，在某种意义下基本对称性是 $\eta\to\eta+2$。

\begin{yizhu}[易错点]$\exp(i\eta\int_DF)$ 里的 $F$ 取实二维形式，对真正的 $U(1)$ 丛有 $\int_DF\in 2\pi\Z$，因此在纯 $U(1)$ 理论中 $\eta\to\eta+1$ 是恒等变换，$\eta$ 的周期本来就是 $1$。正文所谓“基本对称性是 $\eta\to\eta+2$”只有把 $U(1)$ 与 $SU(2)$ 的公共中心 $\Z_2$ 一并计及（等价于把 $U$ 提升为 $SU(2)$ 极大环面的元素）才成立，见式 \eqref{eq:6.13} 前后的讨论。\end{yizhu}


参数 $(\alpha,\beta,\gamma,\eta)$ 取值于 $\mathfrak{K}_0=\R^3\times S^1$。但这并不完全是模型的自然参数空间。我们必须回忆 \eqref{eq:6.3}：模型的 $U(1)$ 对称性可以扩充为 $O(2)$。令 $w$ 为 $O(2)$ 的非连通共轭类中的一个元素。正如我们马上会看到的，$w$ 将扮演 $SU(2)$ 外尔群的非平凡元素的角色。（对更高秩群的余伴随轨道的类似模型，这一事实没有好的类比。这就是我们在此局限于 $SU(2)$ 的原因。）$w$ 通过线性西格玛模型规范群 $U(1)$ 的外自同构作用，并通过改变所有符号作用于耦合参数：
% 式 (6.7)
\setcounter{equation}{6}
\begin{equation}\label{eq:6.7}
  w:(\alpha,\beta,\gamma,\eta)\to
(-\alpha,-\beta,-\gamma,-\eta).
\end{equation}
 这就是本文其余部分中一直来自 $SU(2)$ 外尔群的作用。因此模型的参数空间实际上是 $\mathfrak{K}=\mathfrak{K}_0/\Z_2$。在其不动点处（马上就会描述），$w$ 成为线性西格玛模型的对称性。

\begin{yizhu}[易错点]$w$ 的作用有两层：一是以 $w\,t\,w^{-1}=t^{-1}$（$t\in U(1)$）的方式外自同构线性西格玛模型的规范群 $U(1)$，这是 $U(1)$ 唯一的非平凡外自同构；二是把四个耦合参数整体变号，如式 \eqref{eq:6.7}。$w$ 是 $O(2)$ 中不属于连通分支 $U(1)$ 的元素，故 $w^2\in U(1)$；正是这个 $\Z_2$ 使参数空间成为 $\mathfrak K=\mathfrak K_0/\Z_2$，而 $w$ 只在 $\mathfrak K$ 的不动点处才是模型的对称性。\end{yizhu}


\begin{yizhu}[易错点]约定提醒：$(\alpha,\beta,\gamma)$ 是 $U(1)$ 向量多重态的 Fayet–Iliopoulos 参数（即 $D$-项），$\eta$ 是同一 $U(1)$ 的 $\theta$ 角；这里的 $U(1)$ 是线性西格玛模型的规范群，不是 R-对称性，而 $w$ 通过外自同构作用。$(\alpha,\beta,\gamma,\eta)$ 取值于 $\R^3\times S^1=\mathfrak{K}_0$，而物理参数空间是 $\mathfrak{K}_0/\Z_2$：$w$ 同时改变四个参数的符号。通常的 $\theta$ 周期性给出 $\eta\to\eta+1$；只有把 $U(1)$ 与 $SU(2)$ 的公共 $\Z_2$ 一并考虑（或把 $U$ 提升到 $SU(2)$ 极大环面）时，才看到基本对称性 $\eta\to\eta+2$，这一点在讨论离散电通量 ${\bf e}_0$ 时至关重要（见式~\eqref{eq:6.11}--\eqref{eq:6.13}）。\end{yizhu}

模型在参数空间何处变为奇异，这个问题在 \cite{ref81} 中讨论过。答案是：它恰在点 $P_0:\alpha=\beta=\gamma=\eta=0$ 处奇异；此时由于经典真空的“第二分支”起作用，真空变得不可归一化。

显然，$P_0$ 是 $\Z_2$ 作用于 $\mathfrak{K}_0$ 的两个不动点之一。另一个是 $P_1:\alpha=\beta=\gamma=0$、$\eta=1/2$。那里会发生什么？{}从经典几何的观点看，看不到参数 $\eta$。在 $\alpha=\beta=\gamma=0$ 处的超 Kähler 商是轨形奇点 $\R^4/\Z_2$。与这个轨形（orbifold）相伴有一个可解的共形场论~\cite{ref83}。人们可以猜测，对某个 $\eta$ 值，线性西格玛模型在低能下会约化为这个轨形理论。事实上这确实发生~\cite{ref84}：在 $\eta=1/2$ 处，更确切地说在 $P_1:(\alpha,\beta,\gamma,\eta)=(0,0,0,1/2)$ 处。该点的理论是完全光滑、行为良好的量子场论，而且实际上简单且精确可解。由于 $P_1$ 是 $w$ 的不动点，轨形理论以 $w$ 为对称性。事实上，$w$ 是该轨形的“量子对称性”（它对扭曲扇区中的弦作用为 $-1$，对非扭曲弦作用为 $+1$）。

现在考虑由该理论导出的 $A$-模型或 $B$-模型（选取辛结构或复结构之一），以及参数变动时相应的单值性。事实上，正如本节概述中所解释的，考虑 $A$-模型并不很有趣，因为在当前的局部模型中 $\alpha$ 取值于线性空间，从而排除了有趣单值性的存在。因此我们将专注于某个复结构中的 $B$-模型。由于存在旋转复结构的 $SO(3)$ 对称性，选哪一个都无关紧要。为与本文其余部分的术语一致，我们考虑复结构 $J$ 中的 $B$-模型。

本模型的单值性如同第~\ref{sec:4.5}~节那样，通过固定 $\alpha$ 与 $\gamma$、变化 $\beta$ 与 $\eta$ 而得到。若 $(\alpha,\gamma)$ 非零，则相关参数是模去对称性 $\eta\to \eta+1$ 的 $\beta$ 与 $\eta$。从“外尔变换” $w$ 得不到新东西，因为它不保持 $(\alpha,\gamma)$ 固定。因此在这种情况下，适当的参数空间就是 $\beta-\eta$ 平面对 $\eta\to \eta+1$ 的商，即 $\R^2/\Z=\R\times S^1$。单值群就是 $\Z$。这一答案在第~\ref{sec:4.5}~节中已很熟悉。

照例，最有趣的情形是 $(\alpha,\gamma)=(0,0)$。此时在按 $\Z$ 取商得到 $\R\times S^1\cong\C^*$ 之后，我们还须按 $w:(\beta,\eta)\to (-\beta,-\eta)$ 取商。$w$ 作用在 $\lambda=\exp\left(2\pi(\beta+i\eta)\right)$ 上为 $\lambda\to\lambda^{-1}$。有两个不动点，位于 $\lambda=1$ 与 $\lambda=-1$。$\lambda=1$ 处的不动点就是理论变为奇异的点 $P_0$。我们从参数空间中删去该点，并把剩下的部分记作 $\tilde\C^*$。$\Z_2$ 仍作用于 $\tilde \C^*$，且在 $P_1:\lambda=-1$ 处有一个不动点。我们不想删去这个不动点，因为量子场论在那里定义良好。相反，我们把量子场论所依赖的参数空间看作轨形 $\tilde \C^*/\Z_2$，而相关的单值群是轨形意义下的基本群（fundamental group）$\pi_1^{\rm orb}(\tilde \C^*/\Z_2)$。该群由 $\eura$（绕 $\lambda=1$ 的单值性）与 $\eurt$（绕 $\lambda=-1$ 的单值性）生成，唯一的关系是 $\eurt^2=1$，它表示 $\lambda=-1$ 是一个 $\Z_2$ 轨形点。

\begin{yizhu}[易错点]两处 $\lambda$ 的定义一致：式 \eqref{eq:6.9} 的 $\lambda=\exp(2\pi i(\eta-i\beta))$ 与正文的 $\lambda=\exp(2\pi(\beta+i\eta))$ 逐项相同，因为 $2\pi i\cdot(-i\beta)=2\pi\beta$。$w$ 把 $\lambda$ 送到 $\lambda^{-1}$，在 $\lambda=1$ 与 $\lambda=-1$ 各有一个不动点：前者即理论奇异的 $P_0$（须删去），后者即 $\Z_2$ 轨形点 $P_1$（保留）。故单值群是轨形基本群 $\pi_1^{\rm orb}(\tilde\C^*/\Z_2)=\langle\eura,\eurt\mid\eurt^2=1\rangle$，$\eura$、$\eurt$ 分别是绕 $\lambda=1$、$\lambda=-1$ 的单值性。由于 $\lambda$ 只依赖组合 $\beta+i\eta$，平移 $\eta\to\eta+1$ 把 $\lambda$ 送回自身；而 $w$ 给出的 $\lambda\to\lambda^{-1}$ 并非 $\beta$ 的整数平移，所以 $\tilde\C^*/\Z_2$ 的轨形点只可能落在 $\lambda=\pm1$。\end{yizhu}


局部模型中看到的这一单值性答案，与第~\ref{sec:4.5}~节得到的答案本质上等价，但需要一些说明。首先，考虑群 $SO(3)$ 的极大环面（maximal torus）$\TT$。它由如下形式的元素构成：
% 式 (6.8)
\setcounter{equation}{7}
\begin{equation}\label{eq:6.8}
  U=\left(\begin{matrix}\cos(2\pi\eta)&\sin(2\pi\eta)&0\\ -\sin(2\pi\eta)&
\cos(2\pi\eta)&0\\ 0&0&1\end{matrix}\right).
\end{equation}
若要求紧群 $SO(3)$，则 $\eta$ 为实数；对复李群 $SO(3,\C)=SL(2,\C)/\Z_2$，则应把 $\eta$ 换成其复化 $\eta-i\beta$。在后一情形，我们可以写
% 式 (6.9)
\setcounter{equation}{8}
\begin{equation}\label{eq:6.9}
  U=\left(\begin{matrix}{1\over 2}(\lambda+\lambda^{-1})&{i\over 2}(\lambda-\lambda^{-1})
&0\\ -{i\over 2}(\lambda-\lambda^{-1})&{1\over 2}(\lambda+\lambda^{-1})&0\\ 0&0&1\end{matrix}\right),
\end{equation}
其中 $\lambda=\exp(2\pi i(\eta-i\beta))\in \C^*$。若 $SO(3)$ 中与 $U$ 交换的子群恰为 $\TT$，我们称元素 $U\in\TT$ 为非常正则的（very regular）；若该群的连通分支等于 $\TT$，则称 $U$ 为正则的。这一判据与我们在 $SO(3)$ 还是 $SO(3,\C)$ 中工作无关。对应于 $\lambda=1$ 的元素 $U=1$ 是非正则的。相反，对应于 $\lambda=-1$ 的元素 $U$ 是正则的，但不是非常正则的，因为它与元素
% 式 (6.10)
\setcounter{equation}{9}
\begin{equation}\label{eq:6.10}
  w=\left(\begin{matrix}1&0&0\\ 0&-1&0\\ 0&0&-1\end{matrix}\right)
\end{equation}
交换；该元素通过 $\lambda\to \lambda^{-1}$ 作用，事实上生成 $SO(3)$ 的外尔群。

\begin{yizhu}[含糊点拨]“正则/非常正则”的判据可用中心化子重述：在 $SO(3)$ 中与式 \eqref{eq:6.10} 的 $w$ 交换的元素全体是 $O(2)=\TT\cup w\TT$，其连通分支为 $\TT$。故 $\lambda=-1$ 对应的 $U$（共轭后即 $w$）是正则的；但它与 $w\notin\TT$ 交换，交换群整体不等于 $\TT$，故不是非常正则。而 $\lambda=1$ 的 $U=1$ 与整个 $SO(3)$ 交换，连正则都不是。该判据在 $SO(3)$ 与复化 $SO(3,\C)$ 中相同，因为中心化子的连通分支在复化下不变。\end{yizhu}


对任意 $\lambda\in \tilde\C^*$ 我们都得到一个正则共轭类。然而，$SO(3)$ 中对应于 $\lambda$ 与 $\lambda^{-1}$ 的共轭类在 $w$ 作用下互相共轭。因此 $SO(3)$ 中正则共轭类的模空间是轨形 $\tilde\C^*/\Z_2$；而且如我们所见，它的轨形基本群 $\pi_1^{\rm orb}(\tilde \C^*/\Z_2)$ 也正是 $(\alpha,\gamma)=(0,0)$ 情形下 $B$-模型的单值群。

\begin{yizhu}[易错点]“非常正则”（very regular）与“正则”的区分在这里是关键：正则只要求与 $U$ 交换的群的连通分支为 $\TT$，而非常正则要求整个交换群就是 $\TT$。$\lambda=-1$ 对应的 $U$ 正则但非常正则，因为它还与 $w$ 交换，而 $w$ 正是外尔群的非平凡元素（式 \eqref{eq:6.10}）。另一方面，正文中把 $\eta$ 模 2 等同于把 $U$ 从 $SO(3)$ 提升为 $SU(2)$ 极大环面的元素，这一步把参数空间从 $SO(3)$ 的版本换成 $SL(2,\C)$ 的版本，两个版本的不动点类型因此改变：提升前 $\lambda=1$ 是奇点、$\lambda=-1$ 是轨形点；提升后（见式 \eqref{eq:6.14} 以下）两个不动点都是奇点。核对时不要把这两种提升混为一谈。\end{yizhu}

在第~\ref{sec:4.5}~节中，考虑单值群中对离散电通量（discrete electric flux）${\bf e}_0$ 作用平凡的子群是方便的。因此需在局部模型中找出这个子群。正如我们在第~\ref{sec:4.1}~节解释的，要在 $SU(2)$ 规范场论中定义 ${\bf e}_0$，需允许一个“扭曲”，使规范丛 $E\to M=\Sigma\times C$ 的结构群为 $SO(3)=SU(2)_{\rm ad}$，但限制到 $q\times C$（$q$ 为 $\Sigma$ 中一点）时提升为 $SU(2)$ 丛。在局部模型中，我们不一定有这种到 $SU(2)$ 的提升，因为局部模型的支撑不是 $q\times C$，而是 $D=\Sigma\times p$（$p\in C$ 为某点）。在定义 ${\bf e}_0$ 时，我们计入满足 $\int_Dw_2(E)\not=0$ 的 $SO(3)$-丛 $E\to D$。

这在局部模型中意味着什么？作用在局部模型超多重态上的群不是 $SO(3)$ 而是 $SU(2)$，所以初看之下，人们会以为不可能用 $w_2(E)$ 非零的 $SO(3)$ 丛 $E$ 来扭曲局部模型。然而，局部模型有规范群 $U(1)$；把 \eqref{eq:6.4} 限制到连通分支，我们看到对称群的全局形式不是 $SU(2)\times U(1)$，而是 $(SU(2)\times U(1))/\Z_2$。$w_2(E)\not=0$ 的 $SO(3)$ 丛 $E\to D$ 不能提升为 $SU(2)$ 丛，但可以提升为结构群是 $(SU(2)\times U(1))/\Z_2$ 的丛。这样做时，$U(1)$ 曲率 $F$ 满足
% 式 (6.11)
\setcounter{equation}{10}
\begin{equation}\label{eq:6.11}
  \int_D{F\over 2\pi}={1\over 2}\int_Dw_2(E)~{\rm mod}~\Z.
\end{equation}

\begin{yizhu}[含糊点拨]式 \eqref{eq:6.11} 右端的 $1/2$ 是全部要点：$w_2(E)\not=0$ 的 $SO(3)$ 丛不能提升为 $SU(2)$ 丛，却能提升为结构群 $(SU(2)\times U(1))/\Z_2$ 的丛，此时 $U(1)$ 曲率的整周期被半整数平移，故 $\int_DF/2\pi$ 只在模 $\Z$ 的意义下等于 ${1\over 2}\int_Dw_2(E)$。等价地 $\int_DF=\pi\int_Dw_2(E)+2\pi\Z$，于是平移 $\eta\to\eta+1$ 使被积函数乘以 $e^{i\pi\int_Dw_2}=(-1)^{\int_Dw_2(E)}$（式 \eqref{eq:6.12}），这正是平移离散电通量 ${\bf e}_0$ 的效果。反之若 $w_2(E)=0$，则 $\int_DF$ 是 $2\pi$ 的整数倍，$U(1)$ 丛是普通的，不需要这种扭曲；可见 ${\bf e}_0$ 与 $\eta$ 平移的纠缠只在 $w_2(E)\not=0$ 的扭曲扇区里出现。\end{yizhu}

回顾 \eqref{eq:6.5}，我们看到这意味着 $\eta\to \eta+1$ 的效果是把路径积分的被积函数乘以
% 式 (6.12)
\setcounter{equation}{11}
\begin{equation}\label{eq:6.12}
  (-1)^{\int_Dw_2(E)}.
\end{equation}
这一操作等价于平移 ${\bf e}_0$。

结果是，基本群 $\pi_1^{\rm orb}(\tilde\C^*)$ 是模型的完整单值群，其中包括平移 ${\bf e}_0$ 的变换。若想找出使 ${\bf e}_0$ 保持固定的单值性子群（这将便于与第~\ref{sec:4.5}~节的结果比较），我们应考虑受对称性 $w:(\beta,\eta)\to (-\beta,-\eta)$ 与
% 式 (6.13)
\setcounter{equation}{12}
\begin{equation}\label{eq:6.13}
  \eta\to\eta+2.
\end{equation}
作用的参数对 $(\beta,\eta)$。这些正是保持 ${\bf e}_0$ 固定的对称性。把 $\eta$ 等同于模 2，等价于把 $U$ 提升为 $SU(2)$ 极大环面的元素而不是 $SO(3)$ 的元素。如此提升后，它与
% 式 (6.14)
\setcounter{equation}{13}
\begin{equation}\label{eq:6.14}
  \hat U=\left(\begin{matrix}\exp({i\pi}(\eta-i\beta))& 0\\ 0&\exp(-{i\pi}(\eta-i\beta))
\end{matrix}\right).
\end{equation}
共轭。外尔群仍以两个不动点作用于对 $(\beta,\eta)$，但存在本质差别。现在的不动点是 $P_0:(\beta,\eta)=(0,0)$ 与 $P_1':(\beta,\eta)=(0,1)$。本质差别在于：外尔群的两个不动点现在都是西格玛模型奇异之处（因为其奇点在 $\eta\to\eta+1$ 下不变）。相反，在提升到 $SU(2)$ 之前，一个不动点是西格玛模型的奇点，另一个是轨形点。

\begin{yizhu}[易错点]$\hat U$ 的两个非正则值 $1$、$-1$ 各对应一个不动点：$P_0$ 处 $(\beta,\eta)=(0,0)$ 给 $\hat U=\mathrm{diag}(1,1)=1$；$P_1'$ 处 $(\beta,\eta)=(0,1)$ 给 $\hat U=\mathrm{diag}(e^{i\pi},e^{-i\pi})=-1$。在 $SO(3)$ 中 $1$ 与 $-1$ 是同一个元素，故 $(\beta,\eta)$ 与 $(\beta,\eta+1)$ 描述同一个 $SO(3)$ 元素；把 $\eta$ 模 $2$ 提升到 $SU(2)$ 只是把这个等价类拆成两点，而这两点都落在奇点上。换言之，提升把 $\eta$ 的周期从 $1$ 改为 $2$，从而把 $SO(3)$ 中同一个元素对应的 $P_0$ 与 $P_1'$ 区分开，而这两处的 $\hat U$ 都不可正则。\end{yizhu}


$\hat U$ 的非正则值是 $1$ 与 $-1$，它们恰好对应于不动点 $P_0$ 与 $P_1'$。因此，当我们删去局部模型奇异之处，并按外尔群及 $\eta\to\eta+2$ 取商时，得到的正是 $SL(2,\C)$ 中正则共轭类的模空间。它的基本群称为 $SL(2,\C)$ 的仿射辫群，记作 $B_{\rm aff}(SL(2,\C))$。该群由元素 $\eura$ 与 $\eurb$ 自由生成，我们可以把它们看作绕 $1$ 与 $-1$ 的单值性。它们之间没有任何关系。

\begin{yizhu}[物理图像补充]$B_{\rm aff}(SL(2,\C))$ 是二元自由群：$\Z_2$ 作用 $\lambda\to\lambda^{-1}$ 把绕 $\lambda=1$ 的生成元 $\eura$ 与绕 $\lambda=-1$ 的生成元 $\eurt\eura\eurt$ 配成一对，而 $B_{\rm aff}$ 在 $\pi_1^{\rm orb}(\tilde\C^*/\Z_2)=\langle\eura,\eurt\mid\eurt^2=1\rangle$ 中恰是 $\eurt\mapsto1$ 的核（指数 $2$）。于是 $\eura$ 与 $\eurb=\eurt\eura\eurt$ 之间没有任何关系，特别没有 $\eura\eurb\eura=\eurb\eura\eurb$ 型的三项关系，见式 \eqref{eq:6.15} 及后文关于 $SL(2,\C)$ 是 \eqref{eq:5.16} 与 \eqref{eq:5.17} 之例外的说明。差别只在仿射 Dynkin 图相邻节点的夹角：夹角为 $\pi$ 时三项辫关系自动失效，夹角为 $2\pi/3$ 时才留下关系。\end{yizhu}


局部模型的完整单值群（包括改变 ${\bf e}_0$ 的变换）是一个扩张：
% 式 (6.15)
\setcounter{equation}{14}
\begin{equation}\label{eq:6.15}
  1\to B_{\rm aff}\to \pi_1^{\rm orb}(\tilde\C^*)\to \Z_2\to 1.
\end{equation}
这里 $B_{\rm aff}$ 被一个 2 阶外自同构扩张，该自同构通过 $\eta\to\eta+1$ 作用，并交换在 $SO(3)$ 中等价的点 $\hat U=1$ 与 $\hat U=-1$。我们称这个自同构为 $\eurt$；它满足 $\eurt^2=1$、$\eurt\eura\eurt=\eurb$、$\eurt\eurb\eurt=\eura$。由于 $\eurb$ 与 $\eurt\eura\eurt$ 相同，扩张 $\pi_1^{\rm orb}(\tilde\C^*)$ 由 $\eura$ 与 $\eurt$ 生成，唯一的关系为 $\eurt^2=1$。这就是我们先前给出的完整单值群的结果。

\begin{yizhu}[原文疑误]原文在式 \eqref{eq:6.15} 及本段把完整单值群写成 $\pi_1^{\rm orb}(\tilde\C^*)$，而前文讨论 $\lambda$ 处轨形点时写作 $\pi_1^{\rm orb}(\tilde\C^*/\Z_2)$。后者才与生成元描述 $\langle\eura,\eurt\mid\eurt^2=1\rangle$ 相符：$\tilde\C^*$ 本身同伦于 $S^1$，基本群只是 $\Z$，阶为 $2$ 的 $\eurt$ 来自 $\lambda=-1$ 处的 $\Z_2$ 轨形点。译文照录原文的两处写法，读时按 $\pi_1^{\rm orb}(\tilde\C^*/\Z_2)$ 理解：严格按 $\tilde\C^*$ 读，$\pi_1(\tilde\C^*)=\Z$ 没有阶 $2$ 的元素，式 \eqref{eq:6.15} 的扩张便不成立。\end{yizhu}

所有这些都与第~\ref{sec:4.5}~节的分析一致。单值群中对 ${\bf e}_0$ 作用平凡的子群是仿射辫群，而完整单值群是仿射辫群被规范群 $G$ 的单连通形式的中心所作的扩张。在当前的例子中，$G=SL(2,\C)$，中心为 $\Z_2$。

为补上一处遗留问题，我们或许应当指出：$SL(2,\C)$ 是式 \eqref{eq:5.16} 与 \eqref{eq:5.17} 所给仿射辫群描述的一个例外。例外的原因是：在 $SL(N,\C)$（$N>2$）中，仿射 Dynkin 图的相邻节点对应于夹角为 $2\pi/3$ 的向量，而 $SL(2,\C)$ 并非如此（其夹角为 $\pi$）。为描述 $SL(2,\C)$ 的情形，我们回到仿射外尔群作为扩张 $\Lambda_{\rm cort}\rtimes \Weyl$ 的定义，其中 $\Weyl$ 是外尔群。对 $SU(2)$，$\Lambda_{\rm cort}\cong\Z$，所以 $\Lambda_{\rm cort}$ 中的向量就是一个整数 $n$。外尔群是 $\Z_2$，由 ${\mathfrak{A}}:n\to -n$ 生成。仿射外尔群由 $\eura$ 与 $\eurr:n\to n+1$ 生成。它们满足 ${\eura}^2=1$ 与 $\eura \eurr\eura=\eurr^{-1}$。等价地，$\AffWeyl$ 由 $\eura$ 与 $\eurb=\eura
\eurr$ 生成，满足 $\eura^2=\eurb^2=1$，且没有其他关系。$\eura$ 与 $\eurb$ 是对应于扩充 Dynkin 图两个节点的反射。要得到 $SL(2,\C)$ 的仿射辫群，只需去掉条件 $\eura^2=\eurb^2=1$，于是 $B_{\rm aff}(SL(2,\C))$ 就是以 $\eura$ 与 $\eurb$ 为两个生成元的自由群。

\begin{yizhu}[易错点]仿射 $A_1^{(1)}$ 的两个节点是 $\alpha_1$ 与其仿射对应 $\alpha_0=-\alpha_1$，仿射 Cartan 矩阵的两行分别是 $(2,-2)$ 与 $(-2,2)$，故两向量夹角为 $\pi$；而 $A_n^{(1)}$（$n>2$，即 $SL(N,\C)$，$N>2$）相邻节点的矩阵元为 $-1$，夹角为 $2\pi/3$。三项辫关系 $\eura\eurb\eura=\eurb\eura\eurb$ 只对夹角 $2\pi/3$ 成立，所以 $SL(2,\C)$ 是例外：去掉 $\eura^2=\eurb^2=1$ 之后其仿射辫群退化为自由群。在 $A_1^{(1)}$ 情形 $\eura$、$\eurb$ 都是穿过原点的反射，它们生成的仿射外尔群是无穷二面体群，有限外尔群 $\Z_2$ 只是它的商。\end{yizhu}


\subsection{瞬子与局部奇点}

%% ==== 原文 §6.3：PDF 第 139 页（印刷第 138 页）====
\label{sec:6.3}

在上述讨论中，局部模型在 $\alpha=\beta=\gamma=\eta=0$ 处奇异（而在其他处非奇异）这一事实起了关键作用。

对此 \cite{ref81} 给出了一种解释，用的是“经典真空的第二分支恰在这些参数全为零时变得重要”这一事实。然而，对 $\alpha=\beta=\gamma=\eta=0$ 的重要性还有另一种标准解释。当 $(\alpha,\beta,\gamma)\not=0$ 时，$A_1$ 奇点经历超 Kähler 消解，产生一个拓扑为 $T^*\mathbb{CP}^1$ 的光滑流形。该余切丛的零截面在复结构
% 式 (6.16)
\setcounter{equation}{15}
\begin{equation}\label{eq:6.16}
  {\mathcal I}={(\alpha I+\beta J+\gamma K)\over
\sqrt{\alpha^2+\beta^2+\gamma^2}}.
\end{equation}
 下全纯，在相反的复结构下反全纯。其面积 $A$ 正比于 $\sqrt{\alpha^2+\beta^2+\gamma^2}$。把 $D$ 映入 $\mathbb{CP}^1\subset T^*\mathbb{CP}^1$ 的全纯或反全纯映射，在复结构 $\mathcal I$ 或 $-\mathcal I$ 下是全纯的。我们把它解释为西格玛模型的瞬子。在其他复结构中，不存在紧的全纯或反全纯曲线。由 $D$ 到 $T^*\mathbb{CP}^1$ 零截面的 $d$ 次全纯映射定义的西格玛模型瞬子，对路径积分的贡献正比于 $q^d$，其中 $q=\exp(- kA+2\pi i \eta)$。（这里 $k$ 是依赖于线性西格玛模型规范耦合的常数。）对 $d$ 求和后，瞬子级数在 $q=1$ 处有极点，换言之在 $(\alpha,\beta,\gamma,\eta)$ 全为零时。这个极点反映了西格玛模型的奇点。

\begin{yizhu}[易错点]$q=\exp(-kA+2\pi i\eta)$ 的两项来源不同：$2\pi i\eta$ 来自 $\theta$ 角项 \eqref{eq:6.5} 在 $d$ 次瞬子（$\int_DF=2\pi d$）上的取值，$A$ 是零截面的面积、正比于 $\sqrt{\alpha^2+\beta^2+\gamma^2}$，$k$ 是依赖线性西格玛模型规范耦合的常数。因 $q$ 只依赖 $\eta$ 模 $1$，$q=1$ 当且仅当 $A=0$ 且 $\eta\in\Z$，即四个参数全为零；$\sum_d q^d$ 在此发散，这就是局部模型奇点的来源。\end{yizhu}


以上概述的是二维中的标准分析。我们这里想考虑这一故事在四维中是什么样子。当前讨论所依据的扭曲 ${\mathcal N}=4$ 超杨–米尔斯理论依赖于参数 $t$。根据 \cite{ref8} 第 3.2 节，四维中瞬子方程的类比是条件 $\EUV^+(t)=\EUV^-(t)=\EUV^0=0$，其中
% 式 (6.17)
\setcounter{equation}{16}
\begin{equation}\label{eq:6.17}
  \begin{aligned}\EUV^+(t)&=(F-\phi\wedge\phi+td_A\phi)^+\\
                   \EUV^-(t)&=(F-\phi\wedge\phi-t^{-1}d_A\phi)^-\\
                    \EUV^0&=D\star \phi.\end{aligned}
\end{equation}

此外，根据 \cite{ref8} 的式 (5.28)（或式 (5.7) 与 (5.11)），若特殊化到 $M=\Sigma\times C$，则四维拓扑场论在 $\Sigma$ 上约化为靶空间为 $\MH(G,C)$ 的二维 $A$-模型，其辛结构是
% 式 (6.18)
\setcounter{equation}{17}
\begin{equation}\label{eq:6.18}
  \omega_t={1-t^2\over
1+t^2}\omega_I-{2t\over 1+t^2}\omega_K.
\end{equation}
 的倍数。该 $A$-模型中的瞬子是对应复结构
% 式 (6.19)
\setcounter{equation}{18}
\begin{equation}\label{eq:6.19}
  I_t={1-t^2\over 1+t^2}I-{2t\over 1+t^2}K.
\end{equation}
 的全纯曲线。

现在纳入支撑在 $D=\Sigma_p=\Sigma\times p$（$p$ 为 $C$ 中一点）上的面算符。我们给该面算符配上参数 $(\alpha,\beta,\gamma)$。从局部模型的观点看，我们预期只在复结构 $\mathcal I$ 中有全纯瞬子。局部模型至少应足以描述那些当 $(\alpha,\beta,\gamma)\to 0$ 时作用量趋于零的瞬子。我们称这些为小作用量瞬子（instantons with small action）。四维的小作用量瞬子应对应于复结构为 $I_t$ 的西格玛模型瞬子。而西格玛模型只在复结构 $\mathcal I$ 中有瞬子。因此，我们预期四维规范场论只有在 $I_t=\mathcal I$ 时才有瞬子——至少才有小作用量瞬子。由 \eqref{eq:6.16} 与 \eqref{eq:6.19}，$I_t=\mathcal I$ 的条件是
% 式 (6.20)
\setcounter{equation}{19}
\begin{equation}\label{eq:6.20}
  {\alpha\over\gamma}={t-t^{-1}\over 2},~~\beta=0.
\end{equation}

\begin{yizhu}[含糊点拨]式 \eqref{eq:6.20} 可直接由比对复结构算出：取 $\beta=0$，则式 \eqref{eq:6.16} 给 $\mathcal I=(\alpha I+\gamma K)/\sqrt{\alpha^2+\gamma^2}$，与式 \eqref{eq:6.19} 的 $I_t$ 逐系数相等要求 $\alpha/\sqrt{\alpha^2+\gamma^2}=(1-t^2)/(1+t^2)$ 且 $\gamma/\sqrt{\alpha^2+\gamma^2}=-2t/(1+t^2)$；两式相除即得 $\alpha/\gamma=(1-t^2)/(-2t)=(t-t^{-1})/2$。若 $\beta\not=0$，$I_t$ 没有 $J$ 分量可与之匹配，故必有 $\beta=0$。注意符号：$I_t$ 中 $K$ 的系数为 $-2t/(1+t^2)$，故 $t>0$ 时 $\gamma>0$；又因 $(-\alpha,-\beta,-\gamma)$ 与 $(\alpha,\beta,\gamma)$ 给出同一个 $\mathcal I$，条件其实只定出比值 $\alpha/\gamma$。顺带一提，$\alpha/\gamma=(t-t^{-1})/2$ 在 $t\to-1/t$ 下不变（因 $(-t^{-1}-(-t^{-1})^{-1})/2=(t-t^{-1})/2$），这与该扭曲理论在 $t\to-1/t$ 下的对称性相符，可用来检验式 \eqref{eq:6.20} 的系数是否抄对。\end{yizhu}


我们希望直接在四维中导出这些条件。

首先，从四维观点看，若 $\beta\not=0$，则 $\EUV^0$ 不可能为零。确实，$\beta\not=0$ 会使 $d_A\star\phi$ 沿 $D$ 带有 delta 函数，因此它不可能为零。这是所要的结果之一；剩下的任务是导出 \eqref{eq:6.20} 中的第一个方程。我们将通过推广 \cite{ref8} 第 3.3 节的消没论证来做这件事。

该文献中所谓的消没定理 1（Vanishing Theorem 1）断言：在方程 \eqref{eq:6.17} 的任意解中（无面算符时），庞特里亚金数（Pontryagin number）为零。这可由如下恒等式推出：
% 式 (6.21)
\setcounter{equation}{20}
\begin{equation}\label{eq:6.21}
  \int_M\,\Tr\left(\EUV^+(t)\wedge
\EUV^+(-t^{-1})+\EUV^-(t)\wedge \EUV^-(-t^{-1})\right)=\int_M\Tr
\,F\wedge F.
\end{equation}
 该恒等式的推导依赖于分部积分。方程的解使左端为零，而右端自然是庞特里亚金数的倍数。因此该恒等式意味着庞特里亚金数为零。

\begin{yizhu}[含糊点拨]式 \eqref{eq:6.22} 的最后一项展开后是两块：$4\pi\int_D\Tr\alpha F-2\pi(t-t^{-1})\int_D\Tr\gamma F$，其中第二块正是上一段所说的、定域在 $D$ 上的表面贡献。在 $M=\Sigma\times C$、$D=\Sigma\times p$ 时 $D\cap D=0$ 且 $\mathfrak N=0$，故该项必须为零；对 $G=SU(2)$，$F|_D$ 由 $\mathfrak m$ 刻画，若 $\mathfrak m\not=0$ 就要求 $\alpha-{1\over 2}(t-t^{-1})\gamma=0$，即式 \eqref{eq:6.20}。若 $D\cap D\not=0$（例如 $D$ 有自交），第二项 $4\pi^2(D\cap D)\Tr\alpha^2$ 也进入条件，$\alpha$ 与 $(t-t^{-1})\gamma/2$ 的关系须相应修改；正文只处理 $D=\Sigma\times p$ 的干净情形。\end{yizhu}


存在支撑在 $D\subset M$ 上的面算符时，有几处不同。如同在 \eqref{eq:2.42} 的推导中那样，我们把 \eqref{eq:6.21} 中的积分理解为对 $M$ 中 $D$ 的补集的积分。在这样理解之后，庞特里亚金数就不是 $\int_M\Tr\, F \wedge F$ 的简单倍数，而是带有 \eqref{eq:2.42} 中指示的额外项。其次，在导出 \eqref{eq:6.21} 所需的分部积分中会遇到一项 $(t-t^{-1})\int_M d\left({\rm Tr}\,
\phi\wedge F\right)$。对 $\gamma\not= 0$，尽管它是一个恰当形式的积分，却并不为零，即使 $M$ 紧致也是如此。取而代之的是某种定域在 $D$ 上的表面贡献，其形式为 $-2\pi(t-t^{-1})\int_D\Tr\,\gamma F$。在这样理解之后，\eqref{eq:6.21} 的推广就是
% 式 (6.22)
\setcounter{equation}{21}
\begin{equation}\label{eq:6.22}
  \begin{aligned}\int_M &\,\Tr\left(\EUV^+(t)\wedge
\EUV^+(-t^{-1})+\EUV^-(t)\wedge \EUV^-(-t^{-1})\right)\\ &
 =-8\pi^2\mathfrak{N}+4\pi^2(D\cap D)\Tr\,\alpha^2
+4\pi\int_D\Tr\,\left(\alpha-{t-t^{-1}\over 2}\gamma\right)F.\end{aligned}
\end{equation}


任何超对称方程的解都使左端为零。于是显然，若对小的 $(\alpha,\beta,\gamma)$ 存在这样的解，整数 $\mathfrak{N}$ 必须为零。此外，在我们的应用 $M=\Sigma\times C$、$D=\Sigma\times p$ 中，$D\cap D=0$。因此右端除最后一项外都为零，从而最后一项也必须为零。$G$-丛 $E\to M$ 限制到 $D$ 上时结构群为 $\TT$，并由其曲率 $F$ 的特征类 $\mathfrak{m}$ 在拓扑上刻画。对 $G=SU(2)$（如 \eqref{eq:6.20} 的推导所假定的），$\mathfrak{m}$ 等价于西格玛模型的瞬子数（即映射 $D\to\mathbb{CP}^1$ 的次数）；而西格玛模型在 $(\alpha,\beta,\gamma)$ 趋于零时的奇点，被认为来自 $\mathfrak{m}\not=0$ 的小作用量瞬子。但 \eqref{eq:6.22} 告诉我们，若 $\mathfrak{N}=D\cap D=0$，则 $\mathfrak{m}\not=0$ 的瞬子必须满足 $\alpha-{1\over 2}(t-t^{-1})\gamma=0$。这正是式 \eqref{eq:6.20} 中想要的条件，于是我们从四维观点解释了如下事实：只有当这一条件被满足时，瞬子或超对称场构型才会出现。

\begin{yizhu}[易错点]注意 $\mathfrak{N}$（花体 N）与超对称数 ${\mathcal N}=4$ 中的 $N$ 不是一回事：前者正是庞特里亚金数（决定瞬子数），\eqref{eq:6.22} 右端第一项 $-8\pi^2\mathfrak{N}$ 即来自它；第二项 $4\pi^2(D\cap D)\Tr\,\alpha^2$ 是 $D$ 的自交项（在 $M=\Sigma\times C$、$D=\Sigma\times p$ 时为零）；第三项 $4\pi\int_D\Tr(\alpha-{t-t^{-1}\over 2}\gamma)F$ 是定域在 $D$ 上的表面项，正是它给出条件 \eqref{eq:6.20}。另需注意 $\mathfrak{m}$ 是 $F$ 在 $D$ 上的特征类（对 $G=SU(2)$ 即瞬子数或映射度），与 $\mathfrak{N}$ 含义不同；原文在 \eqref{eq:6.22} 之后把两者并用，核对时勿混淆。\end{yizhu}

同样可以通过把 \cite{ref8} 的消没定理 2 推广到含面算符的情形，导出相同的条件。



\subsection{一些弦论构造}

 %% ==== 原文 §6.4：PDF 第 142 页（印刷第 141 页）====
\label{sec:6.4}

局部模型不保持 $S$-对偶这一断言有点令人困惑。为更好地理解这一点，我们将考虑一个 $S$-对偶显而易见的弦论构造，并看看它为何未能给出局部模型。

如同 \cite{ref18} 那样，我们将考虑 IIB 型超弦理论中一个被广泛研究的情形：一叠 $N$ 张 $D3$-膜给出四维 $U(N)$ 规范场论，它与另一张 $D3$-膜在余维二处相交，后者我们称为 $D3'$。于是所有事情都发生在 IIB 型理论十维时空的某个六维子空间中。我们把这个六维子空间取为三个黎曼面的乘积 $Z=\Sigma\times C\times\tilde C$。$N$ 张 $D3$-膜支撑在 $M=\Sigma\times C\times r$ 上，其中 $r$ 是 $\tilde C$ 中一点；$D3'$-膜支撑在 $\tilde M=\Sigma\times p\times \tilde C$ 上，其中 $p$ 是 $C$ 中一点。$\tilde C$ 既可以是复平面 $\C$，也可以是紧黎曼面。在我们通常的应用中，$C$ 是我们考虑几何朗兰兹纲领的黎曼面，而 $\Sigma$ 是在 $C$ 上紧化所得二维有效场论的时空。也可以把 $M=\Sigma\times C$ 换成更一般的四维流形，眼下这将很方便。

缠绕在 $M$ 上的 $N$ 张 $D3$-膜堆（stack）在 $M$ 上产生具有 ${\mathcal N}=4$ 超对称的 $U(N)$ 规范场论；$D3'$-膜在 $\tilde M$ 上产生同样具有 ${\mathcal N}=4$ 超对称的 $U(1)$ 规范场论。在交点 $M\cap \tilde M=\Sigma\times p\times r$ 上，得到处于 $U(N)$ 的 $N$ 维表示、且在 $U(1)$ 下荷为 1 的超多重态。为尝试与本文所探讨的情形建立联系，人们会希望这些超多重态连同支撑在 $\tilde M$ 上的 $U(1)$ 规范多重态，在低能下等价于靶为 $\mathbb{H}^N/\neg/\neg/U(1)$ 的西格玛模型，其中 $\mathbb{H}^N\cong \C^{2N}$ 由这些超多重态参数化，而 $\mathbb{H}^N/\neg/\neg/U(1)$ 是它对 $U(1)$ 的超 Kähler 商。

\begin{yizhu}[物理图像补充]这一膜构型的取向值得记清：$N$ 张 $D3$-膜沿 $M=\Sigma\times C\times r$ 延展（在 $\tilde C$ 方向局域于一点 $r$），其世界体积上给出四维 $U(N)$ 规范群；单张 $D3'$ 沿 $\tilde M=\Sigma\times p\times\tilde C$ 延展（在 $C$ 方向局域于 $p$），给出 $U(1)$。两者在 $M\cap\tilde M=\Sigma\times p\times r$（即 $\Sigma$ 中的二维曲面）上相交，交会处出现处于 $U(N)$ 的 $N$ 维表示、在 $U(1)$ 下荷为 1 的超多重态；当 $N=2$ 时这正是第~\ref{sec:6.2}~节 $\mathbb{H}^2/\neg/\neg/U(1)$ 线性西格玛模型的场内容，其靶空间在非零 $D$-项下为 $T^*\mathbb{CP}^{N-1}$。$N$ 张膜说到底是 $G=U(N)$、$\mathbb{L}=U(N-1)$ 的局部模型，靶空间为 $T^*(U(N)/U(N-1))=T^*\mathbb{CP}^{N-1}$；但如正文随后指出，该设置因 IIB 的 $SL(2,\Z)$ $S$-对偶而无法真正等价于局部模型。\end{yizhu}

这个超 Kähler 商依赖于超 Kähler 矩映射三个分量所取的常数值。在物理上，这些常数必须来自 $\tilde M$ 上 $U(1)$ 规范场论的 Fayet–Iliopoulos $D$-项；人们希望以某种方式从弦论构造的参数导出它们。对非零 $D$-项，该超 Kähler 商在某个复结构下等价于 $T^*\mathbb{CP}^{N-1}$，在另一个复结构下等价于李代数 ${\mathfrak{gl}}(N)$ 中元素
% 式 (6.23)
\setcounter{equation}{22}
\begin{equation}\label{eq:6.23}
  \delta= \left(\begin{matrix}N-1&&&\\
                 & -1&& \\
                 &&\ddots&\\
                  &&&-1\end{matrix}\right)
\end{equation}
的轨道。这里得到的是第~\ref{sec:6.1}~节所讨论局部模型的一个特例，即 $G=U(N)$、$\mathbb{L}=U(N-1)$ 的情形，于是 $T^*(G/\mathbb{L})=T^*(U(N)/U(N-1))=T^*\mathbb{CP}^{N-1}$。在此情形，参数 $(\alpha,\beta,\gamma)$ 必须各自是 $\delta$ 的实倍数，比例系数由 Fayet–Iliopoulos 参数决定。因此，每个参数只依赖于一个实参数；这一断言在下面解释的 反德西特（Anti de Sitter）构造中仍然成立。

\begin{yizhu}[物理图像补充]式 \eqref{eq:6.23} 的 $\delta$ 满足 $\Tr\delta=(N-1)+(N-1)\cdot(-1)=0$，故 $\delta\in\mathfrak{sl}(N)$。它在 $SU(N)$ 中的稳定化子是 $S(U(N-1)\times U(1))$，故其余伴随轨道就是 $\mathbb{CP}^{N-1}$，其余切丛 $T^*\mathbb{CP}^{N-1}$（复维 $N-1$）既是该轨道的超 Kähler 消解，也是极小幂零轨道闭包的消解。这正是 $G=U(N)$、$\mathbb L=U(N-1)$ 的局部模型的靶空间；参数 $(\alpha,\beta,\gamma)$ 只能取 $\delta$ 的实倍数，故只余三个实参数。\end{yizhu}


这看起来是一条从膜得到本文所讨论情形之例子的可行途径。然而，由于膜构造（如同 IIB 型超弦理论中任何 $D3$-膜构型一样）具有 $SL(2,\Z)$ $S$-对偶，而由 $N$ 个超多重态加 $U(1)$ 规范场构成的模型没有，该膜构型不可能真正等价于把 ${\mathcal N}=4$ 超杨–米尔斯理论耦合到由超多重态与规范场构成的二维系统。失效之处取决于 $\tilde C$ 的选取。

若取 $\tilde C=\C$，则我们得不到想要的那种约化，因为 $\tilde C$ 的无限面积意味着无法用二维方式处理 $D3'$-膜。更根本地，该问题的动力学一般涉及膜复合（brane recombination）：$N$ 张 $D3$-膜中的一张与 $D3'$-膜结合，形成光滑（但形变）的膜。定域地，若 $C$ 由复方程 $x=0$ 定义、$\tilde C$ 由 $y=0$ 定义，则复合可描述为把交点 $xy=0$ 形变为光滑曲线\footnote{对小 $\epsilon$，在以 $x$、$y$ 参数化的空间 $\C^2$ 中，存在一个超对称圆盘，其面积正比于 $\epsilon$，边界位于曲线 $xy=\epsilon$ 上。也可以考虑这个圆盘的多重覆盖。这给出了第~\ref{sec:6.3}~节所讨论小作用量瞬子的一个弦论模型。}
% 式 (6.24)
\setcounter{equation}{23}
\begin{equation}\label{eq:6.24}
  xy=\epsilon.
\end{equation}
于是，低能有效物理无法在不考虑 $\tilde C$ 非紧性的情况下描述，也不会给出本文所考虑的那种面算符。就有利的一面而言，$\tilde C$ 的非紧性意味着从该构造得到的一切都依赖于参数 $(\alpha,\beta,\gamma)$，而这些参数可以用 $\tilde C$ 上无穷远处的行为来定义。例如，$\alpha$ 是 $U(1)$ 规范场在 $\tilde C$ 上无穷远处的和乐，而 $\beta+i\gamma$ 是式 \eqref{eq:6.24} 中的参数 $\epsilon$，它同样可在无穷远处测量（由膜被“弯曲”的方式给出）。

\begin{yizhu}[物理图像补充]以 $x,y$ 记 $\C^2$ 的两个复坐标：$xy=0$ 是两条过原点的复直线之并，而 $xy=\epsilon$（$\epsilon\not=0$）解为 $y=\epsilon/x$，是 $\C^*$ 在该曲面中的图，故曲线光滑且非紧。脚注中面积正比于 $\epsilon$、边界落在该曲线上的超对称圆盘，正是小作用量瞬子的膜论实现：$\epsilon$ 越小面积越小，$\epsilon\to0$ 时圆盘退化，对应 $\beta+i\gamma\to0$。\end{yizhu}


若取 $\tilde C$ 紧致，我们则遇到相反的问题。$\tilde C$ 的紧致性意味着，在低能描述中，$\tilde C$ 上除零模以外的所有场都可被消去。于是我们得到某种描述：$M=\Sigma\times C$ 上的四维场，加上支撑在 $\Sigma\times p$ 上的杂质场。特别地，$D3'$-膜上的 $U(1)$ 规范场在低能下可通过取超 Kähler 商消去。但当 $\tilde C$ 紧致时，描述膜交会处物理的所有变量以及 $\tilde C$ 上的场都是动力学变量（而不是在无穷远处被“冻结”）。因此，没有明显的方式引入对应于 $(\alpha,\beta,\gamma,\eta)$ 的参数。不过，知道从紧致 $\tilde C$ 情形会得到何种面算符，将会是有趣的。

\bigskip\noindent{\it 另一种方案}

对此还有一种替代方案，见 \cite{ref18}。该替代方案使我们既能保持 $S$-对偶，又能看到参数 $(\alpha,\beta,\gamma,\eta)$。它并不等价于局部模型（几乎不可能等价，因为局部模型不是 $S$-对偶的），但它是用膜与弦论描述本文所研究构造之一例的优雅方式。

该替代方案只是利用 AdS/CFT 对偶~\cite{ref85}，把 $D3$-膜替换为涉及这些 $D3$-膜所产生的 IIB 型几何的描述。$D3'$-膜将作为 IIB 型描述的一部分保留下来，从而给出所需的面算符。

为简化，我们取 $M=S^4$，此时最简单相关的 IIB 型时空是 $AdS_5\times S^5$。$D3'$-膜最简单的选取是 $AdS_3\times S^1$，其中 $AdS_3\times S^1$ 以显然的方式嵌入 $AdS_5\times S^5$。$S^1$ 在 $S^5$ 中的嵌入在拓扑上不稳定，但在能量上稳定，事实上给出 $AdS_5$ 上的一个超对称膜。$AdS_3$ 的共形边界是 $D\cong S^2$，它以显然的方式嵌入 $M\cong S^4$（$AdS_3\times S^1$ 的 $S^1$ 因子在共形边界附近事实上收缩为一点）。$D$ 当然是 $M$ 的二维子流形。通过在存在 $D3'$-膜（要求其渐近行为与 $AdS_3\times S^1$ 一致）的情况下作量子化，就在 $M$ 上的 ${\mathcal N}=4$ 超杨–米尔斯理论中定义了支撑在 $D$ 上的超对称面算符。这与用 $AdS_5\times S^5$ 中的弦来描述 $M$ 上的超对称威尔逊算符与 't Hooft 算符~\cite{ref86,ref87} 完全平行。可以对此作修改以考虑更一般的 $M$ 或 $D$，但这对我们将要说的并不重要。

\begin{yizhu}[含糊点拨]欧氏化后 $AdS_3$ 即三维双曲空间 $H^3$，其共形边界是 $S^2$；$D3'$ 的世界体积 $AdS_3\times S^1$ 中 $S^1$ 在 $S^5$ 内嵌入并在共形边界附近收缩，故边界恰为 $D\cong S^2\subset S^4=M$，这正是面算符的二维支撑。这与用 $AdS_2\times S^1\subset AdS_5\times S^5$ 中的弦描述超对称威尔逊算符与 't Hooft 算符~\cite{ref86,ref87} 完全平行：那里 $AdS_2$ 的边界是一维曲线，这里 $AdS_3$ 的边界是二维曲面。\end{yizhu}


这样得到的面算符具有熟悉的参数 $(\alpha,\beta,\gamma,\eta)$。为看出它们，考虑支撑在 $D3'$-膜上、规范群为 $U(1)$ 的 ${\mathcal N}=4$ 超杨–米尔斯理论。在我们的简单例子中，该膜的世界体积是 $X=Y\times S^1$，其中 $Y=AdS_3$。一般地，这给出正确的渐近行为，而描述参数 $(\alpha,\beta,\gamma,\eta)$ 所需的全部不过如此。它们如下产生。参数 $\alpha$ 是 $U(1)$ 规范场绕 $S^1$ 的和乐（在 $Y$ 的共形边界附近测量）。$X$ 上 $U(1)$ 规范场中在绕 $S^1$ 旋转下不变的部分，是 $Y$ 上的一个 $U(1)$ 规范场，它可对偶化为一个角值标量场。该标量场在 $Y$ 边界附近的期望值就是 $\eta$。最后，为找出 $\beta$ 与 $\gamma$，我们使用 $D3'$-膜上规范多重态的“希格斯场” $\varphi$。在庞加莱坐标中，$AdS_3$ 的度规可写成 ${1\over r^2}(dr^2+dx^2+dy^2)$，其中 $0\leq r\leq \infty$，$(x,y)$ 平面是去掉无穷远点的二维球面 $D$。于是 $X=Y\times S^1$ 的度规为 ${1\over
r^2}(dr^2+dx^2+dy^2)+d\theta^2$，其中 $\theta$ 是角度。令 $z=re^{i\theta}$，我们通过要求 $\varphi$ 在 $r=0$（即 $Y$ 的共形边界）处渐近于 $(\beta+i\gamma)(dz/2z)$ 来引入 $\beta$ 与 $\gamma$。可以验证，我们在第~\ref{sec:2.4}~节回顾过的 $U(1)$ 规范场论的基本 $S$-对偶，按预期方式作用于这些参数。

\begin{yizhu}[原文疑误]归一化对照：取 $z=re^{i\theta}$，则 $(\beta+i\gamma)\,dz/z$ 的实部为 $\beta\,dr/r-\gamma\,d\theta$，正是式 \eqref{eq:6.1} 中的 $\phi$。但正文此处要求 $\varphi$ 渐近于 $(\beta+i\gamma)(dz/2z)$，其实部只有式 \eqref{eq:6.1} 的一半（相当于把 $\beta,\gamma$ 各缩小一倍）。这或是 $AdS_3$ 侧希格斯场与四维 $\phi$ 的归一化约定不同，或是原文笔误；译文照录原文，把 $(\alpha,\beta,\gamma)$ 与式 \eqref{eq:6.1} 对照时应以 (6.1) 的归一化为准。\end{yizhu}


\bigskip\noindent{\it 另一变体}


对于该构造的另一变体，我们从六维 $(0,2)$ 共形场论~\cite{ref81} 出发。该理论在 $S^6$ 上的情形可用 $AdS_7\times S^4$ 上的 M 理论~\cite{ref85}（更确切地说，用在具有该渐近行为的时空上的 M 理论）来描述，其中 $S^4$ 上有 $N$ 个单位的通量。$(0,2)$ 模型具有面算符~\cite{ref88,ref89}，可通过考虑 $AdS_7\times S^4$ 中的 $M2$-膜来定义，这些膜的世界体积在无穷远处渐近于二维曲面 $D\subset S^6$。通过把 $AdS_7\times S^4$ 换成具有适当渐近行为的时空，可以把 $S^6$ 换成更一般的六维流形 $Z$。

现在取 $(0,2)$ 模型所表述于其上的六维流形为 $Z=M\times T^2$，其中 $M$ 是四维流形。在 $Z$ 上表述的 $(0,2)$ 模型在低能下等价于 $M$ 上规范群为 $U(N)$ 的 ${\mathcal N}=4$ 超杨–米尔斯理论。$Z$ 上支撑为 $M$ 中一个圆乘以 $T^2=S^1\times S^1$ 中两个因子之一的面算符，给出 $M$ 上的威尔逊算符或 't Hooft 算符。（更一般地，$T^2$ 中任选一个测地圆都给出混合的威尔逊–'t Hooft 算符。）然而，$Z$ 上支撑为 $M$ 中曲面 $D$ 乘以 $T^2$ 中一点 $r$ 的面算符，给出 $M$ 上的面算符，其 $\alpha$ 与 $\eta$ 由 $r$ 的选取决定。$S$-对偶来自 $SL(2,\Z)$ 在 $T^2$ 上的几何作用。该构造与前一个构造密切相关，因为 $T^2$ 上的 M 理论对偶于圆上的 IIB 型理论。

\begin{yizhu}[易错点]$(0,2)$ 理论在 $Z=M\times T^2$ 上紧化给出 $M$ 上的 ${\mathcal N}=4$ 超杨–米尔斯理论，规范群 $U(N)$ 的 $N$ 来自 $S^4$ 上的通量单位数。$T^2$ 上 $SL(2,\Z)$ 的几何作用就是四维理论的 $S$-对偶：$T^2$ 的一个测地圆给出威尔逊算符或 't Hooft 算符，而 $T^2$ 中一点 $r$ 给出面算符，其 $\alpha$ 与 $\eta$ 由 $r$ 的选取决定。$T^2$ 的一个因子缩小时 11 维理论退化为圆上的 IIB 型理论，这正是本变体与前面膜构造的对应。\end{yizhu}

